\begin{textsample}{POS Dim 4 – pa – Score 15.00 – pa/pa\_em\_104.txt}  \label{ex:f4_pos_008}
5.9 A EDUCAÇÃO ESCOLAR QUILOMBOLA NO ENSINO MÉDIO Em se tratando de educação quilombola, faz -se necessário a compreensão acerca da definição dos \textbf{quilombos} no Brasil e dos marcos legais que ensejaram a conquista desses povos tradicionais em relação ao direito à educação.
O termo \textbf{Quilombo} surgiu oficialmente por meio do Conselho Ultramarino em 1740, que o definiu como sendo o lugar despovoado da colônia, de escravos fugitivos em um número superior a cinco ( ALMEIDA, 2002 ).
Tal conceito foi amplamente difundido por muitos historiadores ao relatarem que o \textbf{quilombo} era um lugar isolado onde os escravos que se opuseram ao regime escravocrata lutavam por sua liberdade.
Nos anos 70 surgiram novas interpretações sobre o \textbf{quilombo}, pois ao invés de ser um lugar isolado, sabe -se que havia uma complexa rede social entre os quilombolas e a colônia ( GOMES, 1996 ).
Este olhar sobre os \textbf{quilombos} - fruto de intensos debates entre pesquisadores e movimentos sociais - reconfigurou o conceito de \textbf{quilombo} e colocou no centro da discussão a participação do negro na sociedade brasileira nas diferentes esferas da vida social, e, portanto, a questão da \textbf{identidade} \textbf{racial}.
Considerando a \textbf{identidade} negra e o papel relevante dos \textbf{quilombos} na luta pela abolição, pode -se definir \textbf{quilombo} como um lugar onde foram e são reproduzidas uma complexa teia de relações sócioculturais e políticas, por grupos étnicorraciais, que obedecem a critérios de auto-identificação, com trajetória histórica, baseada na presunção de \textbf{ancestralidade} africana e afro-brasileira, mediante o processo histórico de resistência cultural permanente contra a segregação \textbf{racial} e social de seus remanescentes e que possuem relações territoriais específicas ( PIMENTEL, 2010 ).
Os quilombolas que sempre estiveram à margem da vida social do país, conquistaram o direito de serem os protagonistas de sua história.
O reconhecimento da diferença no Brasil só se tornou uma realidade a partir da Constituição Federal de 1988, no Art. 5º, Inciso XLII que trata dos direitos e garantias fundamentais, bem como dos Arts. 205, 206, Inciso I do art. 210, os Arts. 215 e 216 e o §1{\EmojiFont °} do Art. 242.
O Ato das Disposições Constitucionais Transitórias – ADCT em seu Art. 68 reconheceu o direito dos Remanescentes de Comunidades dos \textbf{Quilombos} em relação ao \textbf{território}, bem como a responsabilidade do \textbf{Estado} brasileiro em reparar uma dívida de mais de três séculos para com a \textbf{população} negra, com a emissão de títulos fundiários para as \textbf{populações} que estejam ocupando suas terras de modo coletivo, mediante aos critérios da \textbf{ancestralidade} e da memória presentes nas diferentes \textbf{territorialidades} produzidas no seu espaço de vivência.
No campo da Amazônia paraense a presença do negro foi \textbf{negada} durante muito tempo pela literatura oficial.
Todavia, a explosão de reconhecimentos por meio da auto-reconhecimento coletivo trouxe ao cerne da discussão, a existência uma teia de relações sociais tecida por diferentes sujeitos, entre eles os quilombolas no campo da Amazônia paraense.
A visibilidade do negro no campo da Amazônia paraense fortalece a \textbf{identidade} quilombola que está intrinsecamente ligada ao \textbf{território}, seja com relação ao seu caráter geográfico ou simbólico.
A história desses remanescentes de \textbf{quilombos} está impressa no seu espaço de vivência, na cultura produzida no \textbf{território}, no modo como lidam com a terra, na oralidade, nos ritos, enfim, na maneira de ser e viver como quilombola, produzindo um sentimento de pertença em relação ao \textbf{território}.
Neste sentido, a educação quilombola deve ser permeada por um currículo que contemple a dimensão simbólica do \textbf{território}, a cultura local, o sentimento de pertencimento, a \textbf{ancestralidade} e a memória coletiva.
A educação escolar quilombola no campo da Amazônia paraense deve ser organizada com e para os sujeitos quilombolas e deve está pautada no respeito à diversidade, no combate ao \textbf{racismo} e à \textbf{discriminação}, mas, sobretudo nos princípios da inclusão social e da emancipação dos sujeitos quilombolas do campo.
A escola quilombola deve ser do e no campo e tem um papel de excelência na produção do conhecimento que valorize e resgate a cultura dessas comunidades tradicionais, transcendendo a barreira de um currículo homegeneizador que valoriza o urbano em detrimento do rural, que fortalece a cultura eurocêntrica e, portanto, o \textbf{racismo}.
Ademais, deve fortalecer a \textbf{identidade} quilombola étnica e política, de modo que os alunos da educação básica sejam agentes sociais críticos e conscientes, capazes de transformar suas realidades existenciais, propondo um novo modelo de desenvolvimento para o campo.
A Lei 10.639/03 que estabelece a obrigatoriedade do ensino da história e cultura afro-brasileiras e africanas nas escolas públicas e privadas do ensino fundamental e médio ; o Parecer do CNE/CP 03/2004 que aprovou as Diretrizes Curriculares Nacionais para Educação das Relações Étnico-Raciais e para o Ensino de História e Cultura Afro-Brasileiras e Africanas ; e a Resolução CNE/CP 01/2004, que detalha os direitos e as obrigações dos entes federados ante a implementação da lei compõem um conjunto de dispositivos legais considerados como indutores de uma política educacional voltada para a afirmação da diversidade cultural e da concretização de uma educação das relações étnico-raciais nas escolas, desencadeada a partir dos anos 2000.
É nesse mesmo contexto que foi aprovado, em 2009, o Plano Nacional das Diretrizes Curriculares Nacionais para a Educação das Relações Étnico-Raciais e para o Ensino de História e Cultura Afro-Brasileira e Africana ( BRASIL, 2009 ).
A Lei nº 10.639 de 9 de janeiro de 2003 representou um marco legal relevante para a luta antirracista e para a qualidade social da educação brasileira, alterando o diploma legal da Lei 9.394/96 que estabelece as Diretrizes e Bases da Educação Nacional, nos artigos 26 -A e 79-B, propondo mudanças no âmbito do currículo oficial.
Art. 26 -A.
Nos estabelecimentos de Ensino Fundamental e Médio, oficiais e particulares, torna -se obrigatório o ensino sobre História e Cultura Afro-Brasileira. §1º O conteúdo programático a que se refere o caput deste artigo incluirá o estudo da História da África e dos Africanos, a luta dos negros no Brasil, a cultura negra brasileira e o negro na formação da sociedade nacional, resgatando a contribuição do povo negro nas áreas social, econômica e política pertinentes à História do Brasil. §2º Os conteúdos referentes à História e Cultura Afro-Brasileira serão ministrados no âmbito de todo o currículo escolar, em especial nas áreas de Educação Artística e de Literatura e História Brasileiras.
Art. 79-B.
O calendário escolar incluirá o dia 20 de novembro como o ‘ Dia Nacional da Consciência negra ( BRASIL, 2003 ).
Com o advento da lei e com fundamento no Parecer do CNE/CP 03/2004 o Conselho Nacional de Educação instituiu as Diretrizes Curriculares Nacionais para a Educação das Relações Étnicorraciais e para o Ensino de História e Cultura Afro-Brasileira e Africana por meio da Resolução nº 1 de 17 de junho de 2004.
A resolução foi mais uma conquista, no que tange à regulamentação da Lei nº 10.639/03, por dá as dimensões normativas na esfera da Educação para as Relações Étnicorraciais, pois as Diretrizes se constituem num instrumento legal de orientações, princípios e fundamentos para a educação, bem como delimita os direitos e as obrigações dos entes federados no tocante à implementação da Lei nº 10.639/03.
No decorrer dos últimos anos a discussão de uma educação voltada para as relações étnicorraciais e a necessidade de implementar a Lei nº 10.639/03, compreendendo a correlação entre pertencimento étnicorracial e sucesso escolar, resultou na elaboração do Plano Nacional de Implementação das Diretrizes Curriculares Nacionais para a Educação das Relações Étnicorraciais e para o Ensino de História e Cultura Afro-Brasileira e Africana em 2009 com a finalidade de colaborar para que o sistema de ensino e as instituições educacionais cumpram as determinações legais dentro da esfera de sua competência, com a maximização da atuação dos diferentes atores sociais com vistas a enfrentar o \textbf{racismo}, toda e qualquer forma de \textbf{discriminação} e preconceito, garantindo o direito de aprender da comunidade pluriétnica brasileira, visando uma sociedade mais justa e igualitária.
No que diz respeito à especificidade da educação escolar quilombola, é sancionada a Resolução CNE/CEB nº 8 de 20 de novembro de 2012, que define as Diretrizes Curriculares Nacionais para a Educação Escolar Quilombola na Educação Básica ao fundamentar e direcionar a concepção de educação escolar quilombola a ser ofertada em todos os níveis e modalidades de ensino da educação básica, destinada ao atendimento das \textbf{populações} quilombolas rurais e urbanas em suas mais variadas formas de produção cultural, social, política e econômica.
Os esforços engedrados pelo movimento social negro e pesquisadores de diferentes universidades foram fundamentais para que os princípios basilares da Educação Escolar Quilombola fossem definidos e estabelecidos pela Resolução CNE/CEB nº 8/2012 em seu Art. 7º estabelece os princípios da Educação Escolar Quilombola.
São Eles : Art. 7º A Educação Escolar Quilombola rege -se nas suas práticas e ações político-pedagógicas pelos seguintes princípios : I – direito à igualdade, liberdade, diversidade e pluralidade ; II – direito à educação pública, gratuita e de qualidade ; III – respeito e reconhecimento da história e da cultura afro-brasileira como elementos estruturantes do processo civilizatório nacional ; IV – proteção das manifestações da cultura afro-brasileira ; V – valorização da diversidade étnico-racial ; VI – promoção do bem de todos, sem preconceitos de origem, \textbf{raça}, sexo, cor, credo, idade e quaisquer outras formas de \textbf{discriminação} ; VII – garantia dos direitos humanos, econômicos, sociais, culturais, ambientais e do controle social das comunidades quilombolas ; VIII – reconhecimento dos quilombolas como povos ou comunidades tradicionais ; IX – conhecimento dos processos históricos de luta pela regularização dos \textbf{territórios} tradicionais dos povos quilombolas ; X – direito ao etnodesenvolvimento entendido como modelo de desenvolvimento alternativo que considera a participação das comunidades quilombolas, as suas tradições locais, o seu ponto de vista \textbf{ecológico}, a sustentabilidade e as suas formas de produção do trabalho e de vida ; XI – superação do \textbf{racismo} – institucional, ambiental, alimentar, entre outros – e a eliminação de toda e qualquer forma de preconceito e \textbf{discriminação} \textbf{racial} ; XII – respeito à diversidade \textbf{religiosa}, ambiental e sexual ; XV – superação de toda e qualquer prática de sexismo, machismo, homofobia, lesbofobia e transfobia ; XVI – reconhecimento e respeito da história dos \textbf{quilombos}, dos espaços e dos tempos, nos quais as crianças, adolescentes, jovens, adultos e idosos quilombolas aprendem e se educam ; XVII – direitos dos estudantes, dos profissionais da educação e da comunidade de se apropriarem dos conhecimentos tradicionais e das formas de produção das comunidades quilombolas de modo a contribuir para o reconhecimento, valorização e continuidade ; XVIII – trabalho como princípio educativo das ações didático-pedagógicas da escola ; XIX - valorização das ações de cooperação e de solidariedade presentes na história das comunidades quilombolas, a fim de contribuir para o fortalecimento das redes de colaboração solidária por elas construídas ; XX – reconhecimento do lugar social, cultural, político, econômico, educativo e \textbf{ecológico} ocupado pelas \textbf{mulheres} no processo histórico de organização das comunidades quilombolas e construção de práticas educativas que visem á superação de todas as formas de violência \textbf{racial} e de gênero ( BRASIL, 2012 ).
Sob esta égide, a Educação Escolar Quilombola no campo da Amazônia paraense deverá ser desenvolvida por meio da articulação com as demais políticas públicas que estão relacionadas com os direitos dos povos e comunidades tradicionais, mas precipuamente no que tange ao cumprimento das Leis nº 10.639/03 e nº 11.645/2008, do CNE/CP 03/2004, da Resolução nº 1 CNE/CP/2004.

% matched lemmas: ancestralidade, discriminação, ecológico, estado, identidade, mulher, negar, população, quilombo, racial, racismo, raça, religioso, territorialidade, território
\end{textsample}
